%=======================================================================
%  Ejercicio 2 — Sobre la población de glóbulos rojos (Tipo 1)
%=======================================================================
\setcounter{ejercicio}{1}   % el enunciado lo numera 2
\begin{ejercicio}[Tipo 1]\label{ej:globulos}
\textbf{Sobre la población de glóbulos rojos.}
En este ejercicio vamos a explorar un modelo para describir el tamaño de la
población de glóbulos rojos en el cuerpo humano. Sea $x_n$ el número de glóbulos
rojos en el cuerpo en el día $n$. El propósito es escribir una función para
describir la cantidad de glóbulos rojos en el día $n+1$. Dicha función debe
quedar descrita en términos de la cantidad de células que se pierden debido a
muerte celular en el día $n$ y en la producción de células el mismo día. Si
$d(x_n)$ representa el número de glóbulos rojos perdidos debido a la muerte
celular en el día $n$ y $p(x_n)$ el número de células rojas ganadas debido a la
producción de la médula ósea en el mismo día, entonces podemos escribir:
\begin{equation}\label{eq:globulos}
  x_{n+1} = x_n - d(x_n) + p(x_n),
\end{equation}
esto es, \emph{el número de glóbulos producidos mañana es el número de glóbulos
rojos de hoy menos aquellos que son destruidos más los que son producidos}. Es
ampliamente aceptado que una fracción constante $c$ de glóbulos rojos es
destruida cada día, por lo que $d(x_n) = c\,x_n$. Donde hay menos información es
en la producción de glóbulos rojos, pero se cree que $p(x_n)$ se comporta como
curvas de Ricker. Básicamente, en baja presencia de glóbulos rojos la médula
ósea produce más mientras que a altas poblaciones de glóbulos produce menos. Dos
posibles formas para $p(x_n)$ son:
\[
  p_1(x) = a\,x\,e^{-bx}, \qquad a>0,\ b>0,
\]
y
\[
  p_2(x) = \frac{b\,\theta^m x}{\theta^m + x^m}, \qquad b>0,\ \theta>0,\ m>0.
\]
\begin{itemize}
  \item Realiza gráficas para $p_1(x)$ y para $p_2(x)$ para diferentes
    parámetros. ¿Qué significados tienen los parámetros en cada función? ¿Qué
    pasa cuando son pequeños y cuando son grandes?
  \item ¿Cómo modifica a largo plazo la selección de parámetros en cada $p(x)$
    a la población total $x_{n+1}$? Compara $p(x)$ contra $d(x)$.
  \item Es sabido que la producción de glóbulos rojos toma varios días. ¿Cómo
    modificarías el modelo para tomar en cuenta ese retraso?
  \item Revisa el artículo \parencite{tetschke2018}. ¿Cuál es la diferencia
    entre el modelo que se propone en ese artículo y el que aquí se presenta?
    Realiza una breve descripción.
  \item En la realidad obtener los parámetros para un modelo puede ser
    complicado. ¿Puedes adaptar algunos parámetros del artículo para este
    modelo?
\end{itemize}
Referencias: \textcite{devries2006}; \textcite{mackey1977};
\textcite{tetschke2018}.
\end{ejercicio}

\begin{solucion}
\begin{enumerate}
  \item Las gráficas las podemos hacer en Python. Con ayuda de la inteligencia artificial
  se facilita el código para graficar las funciones $p_1(x)$ y $p_2(x)$ con diferentes parámetros
  usando widgets en tiempo real. Además, podemos en una misma gráfica mostrar las variaciones
  de los parámetros y ver cómo afectan la forma de las curvas. \textbf{Para el caso de $p_1$, tenemos
  las siguientes curvas:}
  \begin{figure}[H]
    \centering
    \includegraphics[width=\textwidth]{produccion_p1}
    \caption{$p_1(x) = a\,x\,e^{-bx}$: a la izquierda variando $a$ con $b=1$ y
      a la derecha variando $b$ con $a=2$.}
    \label{fig:produccion-p1}
  \end{figure}

  Podemos ver que cuando $a$ es pequeño, el valor máximo de la curva es pequeño, y conforme
  $a$ va creciendo, el valor máximo de la curva también crece. Sin embargo, la curva sigue
  manteniendo una forma similar, cambiando la forma de la curva en la altura del valor máximo
  estirándose de forma vertical. Como $p$ es la cantidad de glóbulos rojos producidos por la 
  médula, el parámetro $a$ podría representar la cantidad de células que la médula
  puede producir por cada glóbulo rojo.

  Por otro lado, cuando $b$ es grande, la médula alcanza su punto máximo con menor cantidad de glóbulos rojos, y conforme $b$ disminuye
  ese punto se recorre más, haciendo que la curva se acerque a 0 más lentamente y crezca verticalmente. Podría
  ser algo similar a la tolerancia de la médula ósea a la cantidad de glóbulos rojos que hay en 
  el cuerpo.

  \textbf{Para el caso de $p_2$, obtenemos las siguientes gráficas:}
  \begin{figure}[H]
    \centering
    \includegraphics[width=\textwidth]{produccion_p2}
    \caption{$p_2(x) = \dfrac{b\,\theta^m x}{\theta^m + x^m}$: a la izquierda
    variando $b$ con $\theta=1$ y $m=4$, al centro variando $\theta$ con $b=2$
    y $m=4$ y a la derecha variando $m$ con $b=2$ y $\theta=1$.}
    \label{fig:produccion-p2}
  \end{figure}

  Similar a la gráfica anterior, el parámetro $b$ para $p_2$ también representa la cantidad de 
  células que la médula puede producir por cada glóbulo rojo. El parámetro $\theta$ cambia
  la escala de la curva, pero mantiene la misma forma por lo que podría representar la tolerancia
  al igual que el parámetro $b$ de $p_1$. Sin embargo, el tercer parámetro sí cambia la forma de la
  gráfica. En los valores $m < 1$, la curva no converge a 0 y muestra un crecimiento infinito. Cuando
  $m = 1$ la curva se acerca asintóticamente a un valor máximo, y cuando $m > 1$ la curva converge a 0. 
  Mientras el valor de $m$ crece en los intervalos $(1, \infty)$, la cantidad de células producidas
  es mayor, pero también se detiene más de golpe.

  \item Cuando graficamos $p(x)$ y $d(x)$, podemos ver que la intersección de ambas curvas nos da un punto de equilibrio donde
  la cantidad de células que se producen es igual a la cantidad de células que se destruyen. \textbf{Para $p_1$ tenemos la siguiente gráfica:}

  \begin{figure}[H]
    \centering
    \includegraphics[width=\textwidth]{largo_plazo_p1_pd.pdf}
    \caption{$p_1(x)$ contra $d(x) = c\,x$ con $c=0.5$: a la izquierda variando $a$ con $b=1$ y a la
      derecha variando $b$ con $a=2$. Los puntos marcan los equilibrios.}
    \label{fig:largo-plazo-p1}
  \end{figure}

  Además, podemos graficar la cantidad de glóbulos rojos en el tiempo usando la ecuación \eqref{eq:globulos} y ver 
  cómo se comporta la población a largo plazo con la variación de los parámetros.

  \begin{figure}[H]
    \centering
    \includegraphics[width=\textwidth]{largo_plazo_p1_poblacion.pdf}
    \caption{Población de glóbulos rojos $x_n$ con $p_1$, $c=0.5$ y $x_0=0.5$: a la izquierda variando $a$
      con $b=1$ y a la derecha variando $b$ con $a=2$.}
    \label{fig:largo-plazo-p1-poblacion}
  \end{figure}

  Cuando variamos $a$ sucede algo interesante. Si la curva de $p_1$ no tiene un punto de equilibrio con $d$,
  la población de glóbulos rojos tiende a 0 eventualmente, mientras que la población de las curvas que sí tienen punto de equilibrio
  tiende a ese valor asintóticamente. De igual manera si recordamos $a$ era la cantidad de células que produce la médula ósea por cada
  glóbulo rojo, por lo que a valores altos, la población de glóbulos rojos oscila de forma más inestable dando brincos más grandes
  en la población, mientras que a valores bajos, la población de glóbulos rojos se mantiene más estable y pegada al valor de equilibrio.

  Por otro lado, al variar la tolerancia (es decir $b$) permitimos que el límite de la población de glóbulos rojos sea más alto o más bajo, pero 
  no afecta la estabilidad de la población. Mantenemos siempre la misma forma de la gráfica solamente variando el punto de equilibrio.
  De manera similar, al variar $c$ variamos la pendiente de la recta y por lo tanto, también modificamos el punto de intersección entre
  $p$ y $d$, por lo que $c$ también es un modificador del equilibrio de la población. En caso de que $a < c$, la población de glóbulos 
  rojos tiende a 0, mientras que si $a > c$, la población tiende a un valor de equilibrio o una oscilación alrededor de ese (como en $a = 40$).

  \textbf{Para $p_2$ tenemos la siguiente gráfica:}
  \begin{figure}[H]
    \centering
    \includegraphics[width=\textwidth]{largo_plazo_p2_pd.pdf}
    \caption{$p_2(x)$ contra $d(x) = c\,x$ con $c=0.5$: a la izquierda variando $b$ con $\theta=1$ y $m=4$,
      al centro variando $\theta$ con $b=2$ y $m=4$, y a la derecha variando $m$ con $b=2$ y $\theta=1$. Los
      puntos marcan los equilibrios.}
    \label{fig:largo-plazo-p2}
  \end{figure}

  Si recordamos, la funcionalidad de los parámetros $b$ y $\theta$ es similar a la de $p_1$, por lo que su efecto en la población de 
  glóbulos rojos es similar al que vimos en la parte anterior. Ahora, para el parámetro $m$, vemos que la reacción de la médula ósea
  es más agresiva, por lo que la población de glóbulos rojos tiende a crecer más rápido, pero llegan a su punto de equilibrio
  de forma más abrupta, cosa que ya habíamos visto previamente en el inciso 1. Con esto, las gráficas de población son las siguientes:

  \begin{figure}[H]
    \centering
    \includegraphics[width=\textwidth]{largo_plazo_p2_poblacion.pdf}
    \caption{Población de glóbulos rojos $x_n$ con $p_2$, $c=0.5$ y $x_0=0.5$: a la izquierda variando $b$
      con $\theta=1$ y $m=4$, al centro variando $\theta$ con $b=2$ y $m=4$, y a la derecha variando $m$
      con $b=2$ y $\theta=1$.}
    \label{fig:largo-plazo-p2-poblacion}
  \end{figure}

  Vemos que para $m$, los valores de $6$ y $10$ generan oscilaciones que no se estabilizan, mientras que para $m=2$ la población de glóbulos 
  rojos tiende a un valor de equilibrio. En particular, cuando $m = 5.3$ la curva de población sí se estabiliza, pero con $m = 5.4$ deja de
  hacerlo, lo que indica que existe un valor especial entre el rango $(5.3, 5.4)$ que hace que la población de glóbulos rojos se estabilice.
  Investigando (principalmente preguntándole a Claude), encontré que ese punto está determinado por la siguiente expresión:

  \begin{desarrollo}[¿De dónde sale la condición para $m$?]
    \textbf{Idea.} El modelo es $x_{n+1} = f(x_n)$ con $f(x) = x - c\,x + p_2(x)$. Cerca del
    equilibrio $x^*$, la distancia $x_n - x^*$ se multiplica cada día por $f'(x^*)$, así que la
    población se estabiliza si
    \[
      |f'(x^*)| < 1 .
    \]

    \textbf{1. Equilibrio.} La producción es igual a la destrucción, $p_2(x^*) = c\,x^*$:
    \[
      \frac{b\,\theta^m}{\theta^m + (x^*)^m} = c
      \qquad\Longrightarrow\qquad
      \theta^m + (x^*)^m = \frac{b\,\theta^m}{c},
      \qquad
      (x^*)^m = \theta^m\left(\frac{b}{c} - 1\right).
    \]

    \textbf{2. Derivada de $p_2$.} Con la regla del cociente, y factorizando:
    \[
      p_2'(x) = b\,\theta^m\,\frac{(\theta^m + x^m) - m\,x^m}{(\theta^m + x^m)^2}
              = \underbrace{\frac{b\,\theta^m}{\theta^m + x^m}}_{=\,c \text{ en } x^*}
                \left(1 - m\,\frac{x^m}{\theta^m + x^m}\right).
    \]

    \textbf{3. Evaluar en $x^*$.} Con las dos igualdades del paso 1,
    \[
      \frac{(x^*)^m}{\theta^m + (x^*)^m} = \frac{\theta^m\,(b/c - 1)}{b\,\theta^m/c} = 1 - \frac{c}{b},
    \]
    así que
    \[
      f'(x^*) = 1 - c + p_2'(x^*)
              = 1 - c + c\left[1 - m\left(1 - \frac{c}{b}\right)\right]
              = \boxed{\,1 - m\,c\left(1 - \frac{c}{b}\right)} .
    \]
    La $\theta$ se cancela: cambia la escala del equilibrio, pero no su estabilidad.

    \textbf{4. Estabilidad.} Cuando hay equilibrio ($b > c$), el término $m\,c\,(1 - c/b)$ es
    positivo, así que $f'(x^*) < 1$ se cumple siempre. La condición que importa es
    $f'(x^*) > -1$, que es la expresión de abajo.
  \end{desarrollo}

  \[
  1 - m\,c\left(1 - \frac{c}{b}\right) > -1 \;\Rightarrow\; m\,c\left(1 - \frac{c}{b}\right) < 2 \;\Rightarrow\; m < \frac{2}{c\,(1 - c/b)} = m_c
  \]

  Donde $m_c$ es ese valor que buscamos. Para los parámetros específicos que estamos usando ($b = 2$ y $c = 0.5$), tenemos que 
  $m_c = \frac{2}{0.5(1 - \frac{0.5}{2})} = \frac{2}{0.5(0.75)} = \frac{2}{0.375} \approx 5.33$. Con esto entonces, encontramos que el punto
  $m_c$ representa el límite de valores de $m$ para que la población de glóbulos rojos sea estable.


  \item La función que estamos usando en este momento es $x_{n+1} = x_n - d(x_n) + p(x_n)$. Ahora, separando cada sección de la función
  tenemos que $x_n$ representa los glóbulos rojos del día $n$, $-d(x_n)$ representa la cantidad de glóbulos rojos que se pierden 
  por muerte celular, y $p(x_n)$ representa la cantidad de glóbulos rojos que se generan. La idea será retrasar la producción de glóbulos
  rojos por unos días, manteniendo la destrucción y cantidad en el día $n$ de glóbulos rojos igual. Ahora, sea $k$ la cantidad de días
  que queremos retrasar la producción de glóbulos rojos, entonces la función se modifica a $x_{n+1} = x_n - d(x_n) + p(x_{n-k})$.

  Algunos cambios que trae este nuevo modelo son que ahora, necesitamos tener $k+1$ condiciones iniciales. También debido a que la cantidad
  producida de glóbulos rojos va retrasada por $k$ días, el equilibrio al que se llega es más inestable debido a la tardía respuesta del
  cuerpo a la cantidad de glóbulos rojos que hay en el cuerpo previo al cálculo para hacer más.

  \item El modelo que se presenta en el artículo tiene ciertas diferencias con el modelo que estamos usando para el ejercicio. Primero, se
  toma en cuenta que el tiempo es continuo, a diferencia del modelo que estamos usando que es discreto. Además, el modelo del artículo
  propone tres ecuaciones para simular cada etapa de las células de forma distinta, y de forma implícita meter el retraso de la producción
  de glóbulos rojos. Además, la idea del modelo que propusieron fue hacer un modelo cuyos parámetros fueran medibles en la vida real,
  en lugar de usar parámetros sin un significado biológico. El modelo que proponen es el siguiente:

  \[
    \begin{aligned}
      \dot x_1 &= \beta\,(X_0 - k_1 x_1) + F_b(x_3)\,x_1, \\
      \dot x_2 &= \beta\,(k_1 x_1 - k_2 x_2), \\
      \dot x_3 &= \beta\,(k_2 x_2 - \alpha\,x_3),
    \end{aligned}
    \qquad
    F_b(x_3) = \gamma\,\frac{\text{Base} - x_3}{\text{Base}}.
  \]

  donde:
  \begin{itemize}
    \item $x_1$ son las células progenitoras que proliferan en respuesta a la eritropoyetina (EPO),
      la hormona que estimula la producción de glóbulos rojos, $x_2$ son las células que siguen
      madurando pero ya no responden a la EPO, y $x_3$ son los glóbulos rojos maduros en la sangre.
    \item $X_0$ es la cantidad constante de células madre que entran a la primera etapa.
    \item $k_1$ y $k_2$ son las tasas de paso de una etapa a la siguiente, y $\alpha = 1/120$ por día 
      es la tasa de muerte de los glóbulos rojos, que viven unos 120 días.
    \item $\beta$ y $\gamma$ son los parámetros propios de cada persona. $\beta$ es qué tan rápida es
      su producción de celulas maduras en general y $\gamma$ qué tan fuerte reacciona a una pérdida de sangre.
    \item $F_b$ representa el efecto de la EPO sin modelarla directamente. Cuando hay menos glóbulos rojos
      que el nivel normal de la persona (\emph{Base}), $F_b > 0$ y las células de $x_1$ proliferan más (se incia producción).
  \end{itemize}

  Además, el modelo del artículo está construido para tener un solo equilibrio estable (\emph{Base}), al que
  la población siempre regresa mientras que en nuestro modelo vimos que la población puede quedarse
  oscilando o incluso comportarse de forma caótica.

  \item Si pensamos en nuestros parámetros con las interpretaciones que les dimos, podríamos utilizar $c$ como
    la tasa de muerte de las celulas, es decir podríamos hacer $c = \frac{1}{120}$. También podríamos usar
    la base $$\text{Base}$$ del artículo para determinar el valor del equilibrio que buscamos, y ajustarlo
    con los parámetros $b$ de $p_1$ y $\theta$ de $p_2$. También se puede interpretar la tolerancia de
    la médula ósea como la fuerza de reacción a la pérdida de sangre, es decir, el parámetro $\gamma$ del artículo.
\end{enumerate}
\end{solucion}
